\subsection*{4.5.}
现在将 (4.4.6) 中的和表示为迹.

对每个 $b\in\mathbb Q$, 定义 $p$ 进 Banach $K$ 代数

\begin{equation}\tag{4.5.1}
L(b)=\text{$K[[X_0,\ldots,X_{n+2}]]$ 中由下述级数组成的子环}
\end{equation}

这些级数可写成

\begin{equation}\tag{4.5.2}
\sum_{dU_0=\sum_{i\geq1}U_i}
B_UX_0^{U_0}X_1^{U_1}\cdots X_{n+2}^{U_{n+2}},
\qquad U=(U_0,\ldots,U_{n+2}),
\end{equation}

且满足

\begin{equation}\tag{4.5.3}
\inf_U\bigl(\operatorname{ord}(B_U)-bU_0\bigr)>-\infty.
\end{equation}

对于下列阶函数 (即范数之对数的相反数), 它是 $p$ 进 Banach 代数

\begin{equation}\tag{4.5.4}
\inf_U\bigl(\operatorname{ord}(B_U)-bU_0\bigr)
=\operatorname{ord}\!\left(\sum_U B_UX^U\right).
\end{equation}

由 (4.4.4) 和 (4.2.1) 得

\[
H(X)\in L\!\left(\frac1{p-1}\right).
\]

引入若干算子. 对每个整数 $k>0$, 用下式定义 $K$ 上 Banach 空间之间的连续算子

\begin{equation}\tag{4.5.5}
\left\{
\begin{aligned}
\psi_k:L(b)&\longrightarrow L(kb),\\
\psi_k\!\left(\sum_U B_UX^U\right)&=\sum_U B_{kU}X^U.
\end{aligned}
\right.
\end{equation}

对每对满足 $b'<b$ 的 $(b,b')$, 有“限制”算子,
它是全连续的 (即有限秩算子的一致极限):

\begin{equation}\tag{4.5.6}
\left\{
\begin{aligned}
L(b)&\xrightarrow{\ \text{限制}\ }L(b'),\\
\sum_U B_UX^U&\longmapsto\sum_U B_UX^U.
\end{aligned}
\right.
\end{equation}

代数 $L(b)$ 的每个元素 $G$ 都给出乘以 $G$ 的连续自同态,

\begin{equation}\tag{4.5.7}
\left\{
\begin{aligned}
G:L(b)&\longrightarrow L(b),\\
\eta&\longmapsto\eta G.
\end{aligned}
\right.
\end{equation}

这些算子由下述引理联系起来, 其证明见 [2] (亦见 [4]).
